\documentclass[10pt,a4paper]{amsart}
\usepackage[margin=19mm]{geometry}
\usepackage{amsmath,amssymb,mathtools}
\usepackage[scr=rsfs,cal=euler]{mathalfa}
\usepackage{xfrac}
\usepackage[unicode,hidelinks,bookmarks=false]{hyperref}
\newcommand{\ie}{i.e.\ }
\newcommand{\cf}{cf.\ }
\newcommand{\resp}{respectively\ }
\newcommand{\Subsection}{\subsection}
\providecommand{\nobreakdash}{\nobreak\hbox{-}}
\setlength{\emergencystretch}{2em}
\multlinegap0pt
\usepackage{bm}
\usepackage{bbm}
\usepackage{dsfont}
\usepackage{xspace}
\usepackage{cancel}
\usepackage{mathtools}
\usepackage{amsthm}
\makeatletter
\@ifundefined{theorem}{\newtheorem{theorem}{Theorem}}{}
\@ifundefined{corollary}{\newtheorem{corollary}[theorem]{Corollary}}{}
\makeatother
\newcommand{\N}{\mathbb{N}}
\newcommand{\PP}{\mathbb{P}}
\newcommand{\R}{\mathbb{R}}
\newcommand{\T}{\mathbb{T}}
\title[Mean-field and fluctuations for hub dynamics in heterogeneou]{Mean-field and fluctuations for hub dynamics in heterogeneous random networks}
\author{Zheng Bian, Jeroen S.W. Lamb, Tiago Pereira}
\date{Source: arXiv:2408.11178v2 (CC BY 4.0)}
\begin{document}
\maketitle

\noindent\emph{Source and licence.} Mean-field and fluctuations for hub dynamics in heterogeneous random networks. Version arXiv:2408.11178v2 (CC BY 4.0), distributed under the Creative Commons Attribution 4.0 International licence, as recorded in the arXiv licence metadata for this exact version and shown on the abstract page https://arxiv.org/abs/2408.11178v2. Full rights notice: licenses/arxiv-2408.11178-CC-BY-4.0.txt.\par\smallskip

\noindent\textit{Editorial scope.} Counted unit: the Corollary whose source label is \texttt{cor:star\_contractions} in the article (source0.tex lines 1380--1387, source label \texttt{cor:star\_contractions}), delivered together with the article's own complete original proof of it (source0.tex lines 1390--1412, one \texttt{proof} environment of 23 lines). No mathematics is written, repaired or re-derived: statement and proof are the article's own text line for line. Required context retained uncounted: eq:hub (lines 711--713); eq:ldn (lines 716--718); eq:reducedhub (lines 731--733); thm:star (lines 751--767); eq:iid\_contractions (lines 1250--1252); eq:iid\_contractions\_coding\_map (lines 1301--1303); thm:typicality\_randomcontractions (lines 1328--1332). Every definition, display and label of those lines is retained unchanged. Omitted from this package and not redistributed: 1--710, 714--715, 719--730, 734--750, 768--1249, 1253--1300, 1304--1327, 1333--1379, 1388--1389, 1413--1725. Nothing inside a retained excerpt is omitted or altered: every retained body is delivered exactly as the article's lines are written, so no omission receipt of any kind accompanies this package. Citations: the retained excerpts carry no citation command at all, so no bibliography entry of the article is transcribed and no reference is redistributed. Reference adaptation: none was needed and none was made; every reference inside the delivered unit resolves inside it, so no unresolved reference or citation is produced. Printed slips kept literal: no printed word, symbol, hypothesis or number of the counted statement or of its proof was altered. Layout and class: the article's own class is \texttt{article} with packages including authblk, fullpage, parskip, tikz, iwona, fontenc, amsmath, amsfonts, hyperref, bbm, mathabx, url, amsthm, graphicx; the wrapper is the lane's pinned \texttt{amsart} editorial wrapper at 10pt on A4, which restates the theorem-like environments the retained text needs and adds the article's own macro definitions that the retained text needs (\texttt{\textbackslash N}, \texttt{\textbackslash PP}, \texttt{\textbackslash R}, \texttt{\textbackslash T}), with extra wrapper packages (bm, bbm, dsfont, xspace, cancel, mathtools). The delivered numbering is the unit's own. Assets: no retained span contains a figure environment, a graphics-inclusion command, a PDF-page inclusion, a table or a TikZ picture; no image file is copied into this package, the package declares no asset dependency and no image file is redistributed with it. Licence: version arXiv:2408.11178v2 of this article is distributed under the Creative Commons Attribution 4.0 International licence, as recorded in the arXiv licence metadata for this exact version and shown on the abstract page https://arxiv.org/abs/2408.11178v2. A full rights notice is delivered at licenses/arxiv-2408.11178-CC-BY-4.0.txt. Characters: every retained excerpt is pure ASCII, so the delivered engine, XeLaTeX with its default Latin Modern text font, reports no missing character.
\begin{equation}\label{eq:hub}
	z^{t+1}= \varphi_1(1,\theta^t\omega,z^t)+ \alpha\cdot\frac{1}{L}\sum_{j=2}^N h(\theta^t \omega,z^t,x_j^t)\mod 1,
\end{equation}

\begin{equation}\label{eq:ldn}
	x_j^{t+1} = \varphi_j(1,\theta^t\omega,x_j^t)  + \alpha \cdot \frac{1}{L} h(\theta^t\omega,x_j^t,z^t) \mod1,~~~~j=2,\cdots,N.
\end{equation}

\begin{equation}\label{eq:reducedhub}
	\varphi_{\alpha,m}(t+1,\omega,z)=\varphi_1(1,\theta^t\omega,\varphi_{\alpha,m}(t,\omega,z)) +\alpha\int_{\T} h(\theta^t\omega,\varphi_{\alpha,m}(t,\omega,z),y)\mathrm{d} m(y).
\end{equation}

\begin{theorem}[Reduction theorem on a star]\label{thm:star}
	Suppose (R1--3) hold for the star network random dynamical system (\ref{eq:hub}-\ref{eq:ldn})  on $N-1$  low-degree nodes and one hub node at coupling strength $\alpha>0$. Let
	initial data $(\omega,x)\in\Omega\times \T^{N}$ be such that the network trajectory $\{x^t=\Phi_{\alpha}(t,\omega,x): t\in\N \}$ admits an $m$-typical shadowing orbit in the low degree coordinates; more precisely, there is $(\omega_{{s}},x_{{s}})\in\Omega\times\T^{N}$ satisfying
	\begin{equation}\tag{Shadowing}
	\sup_{t\in\N}\max_{j=2,\cdots,N}	d_{\T}(x_j^t, \varphi_j(t,\omega_{{s}},x_{{s},j}))\leq \varepsilon_s \left(\alpha (N-1)^{-1} \sup_{t\in\N } \|h(\theta^t\omega,\cdot,\cdot)\|_{C^0}\right); %c_{{s}}\cdot \alpha\cdot  o_1(L)
	%	,~~~~\forall t\in\N ,\forall j=2,\cdots,N;
	\end{equation}
	\begin{equation}\tag{Typicality}
		\frac{1}{T}\sum_{t=0}^{T-1} \delta_{\varphi_2(t,\omega_{{s}},x_{{s},2})} \otimes\cdots\otimes \delta_{\varphi_N(t,\omega_{{s}},x_{{s},N})}\xrightarrow[T\to+\infty]{\text{weak}^*} m^{\otimes (N-1)},
		%~~~~\forall j=2,\cdots,N,
	\end{equation}
	where $d_{\T}$ denotes the distance on the circle, the shadowing precision $\delta_s\mapsto\varepsilon_s (\delta_s )$ is a $\R_+$-valued function converging to 0 as $\delta_s $ tends to 0. Then, for any error tolerance 
	$$\varepsilon\geq  \max\left\{ \alpha (N-1)^{-1/2},4\alpha \sup_{t\in\N }|h(\theta^t \omega,\cdot,\cdot)|_{\mathrm{Lip}} \varepsilon_s \left( \alpha (N-1)^{-1} \sup_{t\in\N } \|h(\theta^t\omega,\cdot,\cdot)\|_{C^0} \right)\right\},$$
	the hub behavior (\ref{eq:hub}) admits $\varepsilon$-reducion to $\varphi_{\alpha,m}$ on initial data $(\omega,x)\in\Omega\times \T^N$ with exceptional asymptotic frequency at most $\rho$ with
	$$ \rho(\varepsilon,\omega)=D(\varepsilon,\omega) \exp(-(N-1)\varepsilon^2 \alpha^{-2} c(\omega)),~~~~D(\varepsilon,\omega)\asymp \varepsilon^{-1} \sup_{t\in\N }\|h(\theta^t\omega,\cdot,\cdot)\|_{C^4},~~c(\omega)=\frac{450}{\sup_{t\in\N } \|h(\theta^t\omega,\cdot,\cdot)\|_{C^4}},$$
	where the constants $D(\varepsilon,\omega)$ and $c(\omega)$ are independent of $N$ and $x$.
\end{theorem}

\begin{equation}\label{eq:iid_contractions}
	\varphi:\N\times\Omega\times M\to M
\end{equation}

\begin{equation}\label{eq:iid_contractions_coding_map}
	\pi:\Omega\to M,~~~~\pi(\omega):=\lim_{n\to+\infty} \varphi(1,\omega,\cdot)\circ\varphi(1,\theta\omega,\cdot)\circ\cdots\circ \varphi(1,\theta^{n-1}\omega,p),~~p\in M,
\end{equation}

\begin{theorem}[Uniform typicality of random orbits]\label{thm:typicality_randomcontractions}
	Consider the iid random iteration $\varphi$ of uniform contractions given in (\ref{eq:iid_contractions}) satisfying (C1-4).	There is a uniform set $\Omega_*\subseteq \Omega$ of full $\PP$-measure such that for every initial condition $x\in M$ and every noise $\omega\in\Omega_*$, the asymptotic behavior of the random orbit $\varphi(t,\omega,x)$ is described by the stationary measure
	$$\frac{1}{T} \sum_{t=0}^{T-1} \delta_{\varphi(t,\omega,x)}\xrightarrow[T\to+\infty]{\text{weak}^*} \pi_{*}\PP,$$
	where $\pi:\Omega\to M$ denotes the coding map given in (\ref{eq:iid_contractions_coding_map}).
\end{theorem}

\begin{corollary}\label{cor:star_contractions}
	Consider star network dynamics (\ref{eq:hub}-\ref{eq:ldn}) satisfying (R1--3), where the node dynamics are given by iid random iteration of contractions $\varphi_i$ in (\ref{eq:iid_contractions}) with uniform contraction rate $\lambda\in(0,1)$ and unique stationary measure $m_0:=\pi_*\PP$ as in Theorem \ref{thm:typicality_randomcontractions}.
	Then, for any error tolerance 
	$$\varepsilon\geq \left\{\alpha L^{-1/2},4\alpha^2 L^{-1} (1-\lambda)^{-1}\sup_{t\in\N} \|h(\theta^t\omega,\cdot,\cdot)\|_{C^0} \sup_{t\in\N } |h(\theta^t\omega,\cdot,\cdot)|_{\mathrm{Lip}}\right\},$$
	the hub behavior admits $\varepsilon$-reduction to $\varphi_{\alpha,m_0}$ defined in (\ref{eq:reducedhub}), for almost every noise realization $\omega$ and any initial condition $x\in\T^N$ with exponentially small exceptional asymptotic frequency at most $\rho$ with
	$$\rho(\varepsilon,\omega)= D(\varepsilon,\omega)\exp(-L\varepsilon^2 \alpha^{-2}c(\omega)),$$
	where $D(\varepsilon,\omega)$ and $c(\omega)$ are the same constants independent of $N$ and $x$ as in Theorem \ref{thm:star}.
\end{corollary}

\begin{proof}
	%In the above dynamical setup, the conditions (R1'--3') in Theorem \ref{thm:loc_star} are satisfied. 
%Consider the uncoupled system $\Phi_{\alpha=0}$.
%\[
%\varphi:\N\times\Omega\times \T^{N-1} \to\T^{N-1},~~~~\varphi(1,\omega,x)_i:=\varphi_i(1,\omega,x_i),~~i=2,\cdots,N.
%\]
By (R1) and (C4), the $\varphi_i$ are independent random uniform contractions on $\T$ with rate $\lambda$. It follows that  the uncoupled system $\Phi_{\alpha=0}$ is a random uniform contraction satisfying (C4) on $\T^{N}$, equipped with the distance $d_{\T^{N}}(x,y)= \max_{i=1,\cdots,N} d_{\T}(x_i,y_i)$, with the same rate $\lambda$. By Theorem \ref{thm:typicality_randomcontractions} applied to $\Phi_{\alpha=0}$ on $(\T^{N},d_{\T^{N}})$, the random orbit $\Phi_{\alpha=0}(t,\omega,x)$ starting from any $x\in\T^N$ and $\omega\in\Omega_*$ asymptotically distributes as the unique stationary measure $m_0^{\otimes N}$. 


To verify the (Typicality) assumption in Theorem \ref{thm:star}, simply take $x_s=x\in\T^N$ and $\omega_s=\omega\in\Omega_*$.
	For the (Shadowing) assumption, we compute
%	in Theorem \ref{thm:star} for the low degree trajectory $\{x_j^t=\Phi_{\alpha}(t,\omega,x)_j:j=2,\cdots,N,t\in\N\}$ and its shadowing random orbit $\varphi_j(t,\omega,x_j)$.
	\begin{align*}
		d_{\T}(x_j^t, \varphi_j(t,\omega,x_j)) \leq & d_{\T}\left(\varphi_j(1,\theta^{t-1}\omega,x_j^{t-1}) + \frac{\alpha}{L}  h(\theta^{t-1}\omega,x_j^{t-1},z^{t-1}), \varphi_j(1,\theta^{t-1}\omega)\circ \varphi_j(t-1,\omega,x_j)\right)\\
		\leq & \frac{\alpha}{L} \sup_{t\in\N} \|h(\theta^t\omega,\cdot,\cdot)\|_{C^0}+ \lambda\cdot d_{\T}(x_j^{t-1}, \varphi_j(t-1,\omega,x_j))\\
		\leq& \alpha  L^{-1} \sup_{t\in\N} \|h(\theta^t\omega,\cdot,\cdot)\|_{C^0}(1+\lambda + \cdot + \lambda^{t-1}) \\
		\leq & \alpha  L^{-1} \sup_{t\in\N} \|h(\theta^t\omega,\cdot,\cdot)\|_{C^0} \frac{1}{1-\lambda},~~~~\forall t\in\N.
	\end{align*}
	We have thus verified for almost every noise realization $\omega$ and for each initial condition $x\in\T^N$ (Shadowing) with %$\varepsilon_s =\alpha\delta\Delta_0^{-1} \sup_{t\in\N} \|h(\theta^t\omega,\cdot,\cdot)\|_{C^0}$ the 
	shadowing precision $\varepsilon_s =  \frac{1}{1-\lambda} \alpha L^{-1} \sup_{t\in\N} \|h(\theta^t\omega,\cdot,\cdot)\|_{C^0}$. 
	%shadowing constant $c_{{s}}= \sup_{t\in\N} \|h(\theta^t\omega,\cdot,\cdot)\|_{C^0} (1-\lambda)^{-1}$. 
	The corollary follows from Theorem \ref{thm:star}.
\end{proof}

\end{document}
